\documentclass[11pt]{article}
\usepackage[a4paper,margin=27mm]{geometry}
\usepackage[T1]{fontenc}
\usepackage{lmodern}
\usepackage{microtype}
\usepackage{amsmath,amssymb,amsthm,mathtools}
\usepackage{booktabs,longtable,array}
\usepackage{xcolor}
\usepackage{enumitem}
\usepackage{hyperref}
\hypersetup{colorlinks=true,linkcolor=blue!50!black,urlcolor=blue!60!black,
  citecolor=blue!50!black,pdftitle={An explicit prime number theorem in intuitionistic primitive recursive arithmetic}}
\setlength{\emergencystretch}{3em}
\setlist{itemsep=3pt,topsep=6pt}
\newtheorem{theorem}{Theorem}[section]
\newtheorem{lemma}[theorem]{Lemma}
\newtheorem{proposition}[theorem]{Proposition}
\newtheorem{corollary}[theorem]{Corollary}
\theoremstyle{definition}
\newtheorem{definition}[theorem]{Definition}
\theoremstyle{remark}
\newtheorem{remark}[theorem]{Remark}
\newcommand{\ipra}{\mathsf{iPRA}}
\newcommand{\N}{\mathbb N}
\newcommand{\Q}{\mathbb Q}
\newcommand{\abs}[1]{\left|#1\right|}
\newcommand{\ceil}[1]{\left\lceil#1\right\rceil}
\newcommand{\floor}[1]{\left\lfloor#1\right\rfloor}
\newcommand{\dist}{\operatorname{dist}}
\newcommand{\dotminus}{\dot{-}}
\title{An explicit prime number theorem in\\
intuitionistic primitive recursive arithmetic}
\author{Prepared for Vladimir Reshetnikov}
\date{November 14, 2026}

\begin{document}
\maketitle
\begin{abstract}
We give a finite, quantitative elementary proof of the prime number theorem
and explain its interpretation in first-order intuitionistic primitive
recursive arithmetic, with induction restricted to quantifier-free formulas.
The actual arithmetic conclusion uses no real numbers: it compares the
prime-counting function with the rational number $H_{n^2}-H_n$, and clears
denominators with $(n^2)!$. An explicit double-exponential convergence modulus
is supplied. The central argument combines a finite harmonic substitute for
Euler's constant, Selberg's symmetry identity, and a contraction established
by a finite geometric mesh with rational comparisons. We give the full
human-readable argument, including the bounded induction and numerical
enclosure bookkeeping needed for the weak constructive setting. The
accompanying Lean files specify the arithmetic statement and the required
deeply embedded proof endpoint; they are an uncompiled formalization outline,
not a completed Lean proof.
\end{abstract}

\paragraph{Reading guide.}
The main arithmetic theorem is stated in Section~\ref{sec:statement}.
The estimates and the contraction are proved in full in the middle sections.
Section~\ref{sec:arithmetic} explains why the displayed analytic notation
denotes finite arithmetic calculations and why the induction is permitted in
$\ipra$. The final section gives a concrete Lean implementation route and
states exactly which files remain to be formalized.

\newpage
\tableofcontents
\newpage

\section{The exact theorem and the theory}\label{sec:statement}

\subsection{An entirely arithmetical statement}

Let $\pi(n)$ count the primes at most $n$, with primality decided by trial
division through the bounded interval $2\le d<n$. Put
\[
 F(n)=(n^2)!,\qquad
 A(n)=\sum_{j=n+1}^{n^2}\frac{F(n)}{j},\qquad
 L(n)=\frac{A(n)}{F(n)}=H_{n^2}-H_n.
\]
Every division in $A(n)$ is exact when $n\ge1$. The empty sum is zero.
The symbol $F$ here is a denominator; it is unrelated to any analytic
function. For natural numbers, write
\[
 \dist(a,b)=(a\mathbin{\dotminus}b)+(b\mathbin{\dotminus}a),
\]
where $\dotminus$ denotes truncated subtraction. Thus $\dist(a,b)=|a-b|$.

For $k\ge1$, define
\begin{equation}\label{eq:modulus}
 q=8k,\qquad D=2^{20}q^3,\qquad
 N(k)=2^{\,2^{D^2+2}},
\end{equation}
and set $N(0)=2$.

\begin{theorem}[Arithmetic prime number theorem]\label{thm:main}
The following sentence is provable in first-order $\ipra$:
\begin{equation}\label{eq:arith-pnt}
 \forall k\,\forall n\;
 \bigl(k\ge1\;\wedge\; n\ge N(k)\bigr)
 \ \longrightarrow\ 
 k\,\dist\bigl(\pi(n)A(n),nF(n)\bigr)\le nF(n).
\end{equation}
All functions occurring in this formula are explicitly primitive recursive.
Its matrix, including its antecedent, can be expressed as one Boolean
primitive-recursive function equalling $1$.
\end{theorem}

The modulus is intentionally coarse. It is a closed elementary function;
we use no unbounded search to obtain a threshold. The factorial denominator
is similarly chosen for transparency, not numerical efficiency. Neither
the proof nor the companion programs attempt to evaluate $N(1)$ in full.

For $n\ge1$ we will prove
\begin{equation}\label{eq:log-surrogate}
 0\le \log n-L(n)\le\frac1n.
\end{equation}
Consequently \eqref{eq:arith-pnt}, applied with $2k$ and with
$n\ge\max\{N(2k),2k,2\}$, implies
\[
 \abs{\frac{\pi(n)\log n}{n}-1}
 \le\frac1{2k}+\frac1n\le\frac1k.
\]
This is the usual prime number theorem. Conversely, the same error estimate
shows that the rational formulation expresses the same asymptotic assertion.
The explicit modulus is additional information.

\subsection{Which intuitionistic primitive recursive arithmetic?}

We use first-order intuitionistic logic with equality and a symbol for every
finite primitive-recursive construction: zero, successor, projections,
composition, and primitive recursion. Different construction codes may name
extensionally equal functions without being identified syntactically.
The nonlogical axioms are the defining equations for these constructions,
$Sx\ne0$, injectivity of $S$, and the induction schema
\begin{equation}\label{eq:qfind}
 \varphi(0,\vec z)\to
 \bigl(\forall u\,[\varphi(u,\vec z)\to\varphi(Su,\vec z)]\bigr)
 \to\forall u\,\varphi(u,\vec z),
\end{equation}
where $\varphi$ is quantifier-free. Order and all other arithmetic notation
are definable by primitive-recursive characteristic functions.

The logical rules are the intuitionistic introduction and elimination rules
for $\wedge,\vee,\to,\forall,\exists$, elimination of falsity, and equality
reflexivity and substitution. Universal introduction and existential
elimination have their usual fresh-variable conditions; implication
introduction discharges its hypothesis. General excluded middle and
double-negation elimination are not rules. Equality, order, primality and
every other bounded arithmetic test used below are decidable by their
explicit Boolean algorithms. Case distinctions about these tests are
therefore intuitionistically valid.

This is exactly the presentation called \texttt{PRAAx} in the inspected
\texttt{ConciseIPRA/PR.lean} source, whose logic is the intuitionistic
deduction system in \texttt{Deduction.lean}; the name \texttt{PRAAx} should
not be mistaken for classical PRA. The repository provenance and the
formalization endpoint are recorded later.

\subsection{What must not be inferred just from a classical proof}

The usual assertion
$\forall k\,\exists b\,\forall n\ge b\;\text{error}(k,n)$ is not made
intuitionistically provable merely by calling its proof elementary.
Nor does $\Pi_2$ conservativity of $I\Sigma_1$ over classical PRA settle
this $\Pi_3$ assertion. Avigad explicitly identifies this problem with the
conservativity claim in Sudac's paper~\cite{avigad,sudac}. The classical
weak-arithmetic result of Cornaros and Dimitracopoulos~\cite{cd} is valuable
background, but we do not use it as a black-box replacement for the proof.

Here the threshold is a displayed function. More importantly, the proof
of its correctness uses bounded arithmetic calculations and quantifier-free
induction directly. We do not invoke a conservation theorem, Markov's
principle, a least real zero, a limsup, or a compactness argument.

\section{Finite logarithms and the meaning of the notation}\label{sec:log}

\subsection{Rational logarithm intervals}

For positive integers $a\ge b$ and $M\ge1$, define
\[
 \lambda_M(a/b)=H_{Ma}-H_{Mb}.
\]
All of these are rational numbers obtained by bounded sums. The elementary
harmonic comparison gives
\begin{equation}\label{eq:harmonic-log}
 0\le\log(a/b)-\lambda_M(a/b)
 \le \frac1{Mb}-\frac1{Ma}\le\frac1M.
\end{equation}
For a rational below $1$, use the negative of the interval for its reciprocal.
Reducing the fraction is optional. All endpoints, including their common
denominators, are primitive recursive.

For completeness, write $h_j=H_j-\log j$. The rectangle inequality
\[
 \frac1{j+1}\le\log\left(1+\frac1j\right)\le\frac1j
\]
implies
\[
 0\le h_j-h_{j+1}\le\frac1j-\frac1{j+1}.
\]
Summation from $Mb$ to $Ma-1$ proves \eqref{eq:harmonic-log}.
Taking $M=n$, $a=n$, $b=1$ proves \eqref{eq:log-surrogate}.
No limit of $h_j$ is needed.

There is also a formulation entirely inside rational arithmetic.
$\lambda_M(a/b)$ is the right rectangle sum for $1/t$ on $[b,a]$
with mesh $1/M$; its upper rectangle sum is
$\lambda_M(a/b)+1/(Mb)-1/(Ma)$. Refinement nests these intervals.
For two unrelated mesh sizes, use their product as a common refinement.
Thus their consistency and width are bounded rational inequalities,
proved without presupposing a real logarithm. Additivity is witnessed by
the exact telescoping identity
\[
 H_{Mac}-H_{Mbd}
 =(H_{Mad}-H_{Mbd})+(H_{Mac}-H_{Mad}).
\]
The two summands are logarithm approximants at the refined mesh sizes
$Md$ and $Ma$. Comparing them with the original intervals gives the
stated finite error form of logarithm additivity.

\subsection{Finite rectangle calculus}

Throughout the argument, integrals over a bounded interval are convenient
notation for finite calculations with these logarithm intervals. Here is
the concrete justification. For positive rational endpoints $a<b$ and a
partition $a=t_0<\cdots<t_m=b$, the lower and upper rectangle sums for
$1/t$ differ by at most
\[
 \max_i(t_i-t_{i-1})\left(\frac1a-\frac1b\right).
\]
This follows by bounding each cell with its endpoint values and telescoping.
Partitions can be concatenated or rescaled using exact rational arithmetic.
Their refinement gives additivity and the identity
$\log(xy)=\log x+\log y$, with an explicit enclosure error at every finite
precision. This is also a direct construction of the logarithm notation
without assuming a real-number domain in $\ipra$.

The only other calculus rules used below are finite summation by parts,
rectangle comparison, and antiderivatives of powers and products of
logarithms. They can be obtained by telescoping on a rational mesh.
For example, on $a\le u\le v\le b$,
\[
 \frac{v-u}{v}\le\log v-\log u\le\frac{v-u}{u},
\]
so replacing the logarithmic increment by $(v-u)/u$ has error at most
$(v-u)^2/a^2$. The algebraic identity
$z^2-w^2=2w(z-w)+(z-w)^2$ then gives the rule for $(\log u)^2$ with an
explicit quadratic cell error. The product identity
\[
 f(v)g(v)-f(u)g(u)
 = f(u)\bigl(g(v)-g(u)\bigr)
  +g(u)\bigl(f(v)-f(u)\bigr)
  +\bigl(f(v)-f(u)\bigr)\bigl(g(v)-g(u)\bigr)
\]
gives the remaining product and integration-by-parts rules. Summing the cell
errors costs at most a constant times the mesh width. Rational enclosures
for square and fourth roots are obtained by bounded bisection; their
properties also follow from polynomial identities. Thus the appearances of
$\sqrt{x}$ and $x^{1/4}$ below require no root-existence axiom.

For a step function $B(u)=\sum_{n\le u}b_n$, the identity
\begin{equation}\label{eq:abel}
 \sum_{n\le x}b_n g(n)
 =B(x)g(x)-\int_1^x B(u)g'(u)\,du
\end{equation}
is simply telescoping on the finitely many intervals between integers.
When $g(1)\ne0$, the convention $B(1^-)=0$ is understood. There is no
improper integral in this paper. The later error budget makes explicit why
these finite translations are sufficient for the arithmetic conclusion.

\section{Finite Selberg estimates with explicit constants}
\label{sb:section}

We now prove all the number-theoretic estimates needed for the contraction
argument.  The constants are deliberately generous.  Their purpose is to
make every error bound visible and to avoid any appeal to an unspecified
asymptotic estimate.

For a positive integer $n$, define
\[
 \Lambda(n)=
 \begin{cases}
  \log p,&n=p^j\text{ for a prime }p\text{ and }j\ge1,\\
  0,&\text{otherwise},
 \end{cases}
 \qquad
 \psi(x)=\sum_{n\le x}\Lambda(n),
 \qquad R(x)=\psi(x)-x.
\]
Thus $\Lambda(1)=0$, and $\psi(x)=0$ for $0\le x<2$.  A summation
with a real upper endpoint means a finite sum over the positive integers
below that endpoint.  All integrals below have bounded endpoints.  They
can be evaluated or approximated by the finite arithmetic procedures
specified in the arithmetization part of the paper.

\subsection{Harmonic estimates without an infinite constant}
\label{sb:harmonic}

Write
\[
 H_m=\sum_{j=1}^m\frac1j,\qquad
 h_m=H_m-\log m,\qquad
 P_i(y)=\sum_{n\le y}(\log n)^i\quad(i=1,2).
\]
We use the elementary inequalities
\begin{equation}
 \frac{z-1}{z}\le\log z\le z-1\qquad(z\ge1).
 \label{sb:log-basic}
\end{equation}
They follow by comparing $1/t$ with its endpoint values on $[1,z]$.
In particular $1/2\le\log2\le1$.

\begin{lemma}
\label{sb:harmonic-finite}
For integers $1\le m\le M$,
\[
 0\le h_m-h_M\le\frac1m-\frac1M,
 \qquad 0<h_M\le1.
\]
For every $y\ge1$,
\begin{align}
 |P_1(y)-(y\log y-y)|&\le\log y+1\le y,
 \label{sb:log-sum}\\
 |P_2(y)-y((\log y)^2-2\log y+2)|
   &\le(\log y)^2+2.
 \label{sb:log-square-sum}
\end{align}
\end{lemma}

\begin{proof}
The difference of successive $h$'s is
\[
 h_n-h_{n+1}=\log(1+1/n)-\frac1{n+1}.
\]
By \eqref{sb:log-basic}, this lies between $0$ and
$1/n-1/(n+1)$.  A finite telescoping sum proves the first assertion.
The rectangle bounds
\[
 \log(M+1)\le H_M\le1+\log M
\]
give the second assertion.

If $f$ is either $\log t$ or $(\log t)^2$, and $N=\lfloor y\rfloor$,
monotonicity on $[1,y]$ gives
\[
 0\le\sum_{n=1}^N f(n)-\int_1^N f(t)\,dt\le f(N).
\]
The extra integral from $N$ to $y$ lies between $0$ and $f(y)$.
Consequently the distance between the sum and $\int_1^y f(t)\,dt$
is at most $f(y)$.  Now use
\[
 \int_1^y\log t\,dt=y\log y-y+1,
\]
\[
 \int_1^y(\log t)^2\,dt
    =y((\log y)^2-2\log y+2)-2.
\]
Finally, $\log y+1\le y$ is another instance of
\eqref{sb:log-basic}.
\end{proof}

\subsection{A Chebyshev bound and the signed logarithmic mean}
\label{sb:chebyshev}

\begin{lemma}
\label{sb:basic-psi}
For $x\ge1$,
\begin{equation}
 0\le\psi(x)\le4x,\qquad |R(x)|\le4x,
 \label{sb:chebyshev-bound}
\end{equation}
and
\begin{equation}
 \left|\sum_{n\le x}\frac{\Lambda(n)}n-\log x\right|\le6.
 \label{sb:mertens-bound}
\end{equation}
For $1\le a\le b$,
\begin{equation}
 \left|\int_a^b\frac{R(u)}{u^2}\,du\right|\le20.
 \label{sb:signed-integral}
\end{equation}
\end{lemma}

\begin{proof}
Counting the multiples of each prime power in a factorial gives the exact
identity
\[
 \log\binom{2n}{n}
  =\sum_{d\le2n}\Lambda(d)
       \left(\left\lfloor\frac{2n}{d}\right\rfloor
              -2\left\lfloor\frac nd\right\rfloor\right).
\]
Every parenthesis is either $0$ or $1$, and it is $1$ when $n<d\le2n$.
The binomial theorem gives $\binom{2n}{n}\le2^{2n}$.  Therefore
\[
 \psi(2n)-\psi(n)\le2n\log2.
\]
Apply this for $n=1,2,\ldots,2^{k-1}$ along the dyadic sequence and
telescope:
\[
 \psi(2^k)\le(2^{k+1}-2)\log2<2^{k+1}.
\]
For $x>1$, the least dyadic integer $2^k\ge x$ is less than $2x$.
Monotonicity gives $\psi(x)\le4x$.  The case $x=1$ is immediate.
Since $R(x)=\psi(x)-x$, this also gives the stated bound on $R$.
For later use, $|R(x)|\le4x$ holds for $0\le x<1$ as well, because
there $R(x)=-x$.

The prime factorization of $n$ gives
\[
 \sum_{d\mid n}\Lambda(d)=\log n.
\]
Summing over $n\le x$ and reversing two finite sums yields
\[
 P_1(x)=\sum_{d\le x}\Lambda(d)
                      \left\lfloor\frac xd\right\rfloor.
\]
Put $A_1(x)=\sum_{d\le x}\Lambda(d)/d$.  The floor errors satisfy
\[
 |xA_1(x)-P_1(x)|\le\psi(x)\le4x.
\]
Together with \eqref{sb:log-sum}, this gives
\[
 |A_1(x)-\log x|
 \le4+\left|-1+\frac{P_1(x)-(x\log x-x)}x\right|
 \le6.
\]

The remaining identity can be checked by integrating each of the finitely
many terms $\Lambda(n)\mathbf1_{[n,x]}(t)$:
\[
 \int_1^x\frac{\psi(t)}{t^2}\,dt
     =\sum_{n\le x}\Lambda(n)\left(\frac1n-\frac1x\right)
     =A_1(x)-\frac{\psi(x)}x.
\]
Subtracting $\int_1^x dt/t=\log x$ gives
\[
 \int_1^x\frac{R(t)}{t^2}\,dt
     =A_1(x)-\frac{\psi(x)}x-\log x.
\]
Its absolute value is at most $10$.  Subtract the expressions at $b$
and $a$ to obtain \eqref{sb:signed-integral}.
\end{proof}

\subsection{A finite hyperbola calculation}
\label{sb:hyperbola}

Let $\tau(n)$ be the number of positive divisors of $n$, and put
$D(y)=\sum_{n\le y}\tau(n)$.  Fix an outer endpoint $x\ge1$ and set
\begin{equation}
 M=\lceil x\rceil,\qquad c=H_M-\log M.
 \label{sb:finite-c}
\end{equation}
The same parameter $c$ will be used for every $1\le y\le x$.
It is a finite expression; no existence or convergence assertion about
Euler's constant is being used.

\begin{lemma}
\label{sb:divisor-approx}
For $1\le y\le x$,
\begin{equation}
 |D(y)-y\log y-(2c-1)y|\le12\sqrt y.
 \label{sb:divisor-bound}
\end{equation}
Furthermore, the arithmetic function
\begin{equation}
 f(n)=\tau(n)\log n-2c\tau(n)+4c^2
 \label{sb:comparison-function}
\end{equation}
satisfies
\begin{equation}
 \left|\sum_{n\le y}\bigl(f(n)-(\log n)^2\bigr)\right|
       \le80\sqrt y(\log y+1).
 \label{sb:comparison-error}
\end{equation}
\end{lemma}

\begin{proof}
Let $m=\lfloor\sqrt y\rfloor$.  Counting pairs of positive integers
whose product is at most $y$, according to which coordinate is at most
$m$, gives
\[
 D(y)=2\sum_{n\le m}\left\lfloor\frac yn\right\rfloor-m^2
     =2yH_m-m^2+e,\qquad |e|\le2m.
\]
By Lemma~\ref{sb:harmonic-finite},
\[
 H_m=\log m+c+\eta_m,\qquad 0\le\eta_m\le1/m.
\]
Thus
\[
 D(y)-y\log y-(2c-1)y
   =-y\log(y/m^2)+(y-m^2)+2y\eta_m+e.
\]
Because $m\le\sqrt y<m+1\le2m$, we have
\[
 y\log(y/m^2)\le\frac{2y}{m}\le4\sqrt y,
 \qquad 0\le y-m^2\le2\sqrt y,
\]
\[
 2y\eta_m\le4\sqrt y,\qquad |e|\le2\sqrt y.
\]
Adding these bounds proves \eqref{sb:divisor-bound}.

For $T(y)=\sum_{n\le y}\tau(n)\log n$, finite partial summation gives
\[
 T(y)=D(y)\log y-\int_1^y\frac{D(t)}t\,dt.
\]
Write $D(t)=t\log t+(2c-1)t+E(t)$, where
$|E(t)|\le12\sqrt t$.  Substitution and elementary integration yield
\[
 T(y)=y\bigl((\log y)^2+(2c-2)\log y+2-2c\bigr)+E_T(y),
\]
where
\[
 |E_T(y)|\le2+12\sqrt y\log y+24(\sqrt y-1)
                 \le30\sqrt y(\log y+1).
\]
Here the constant term before the error is $2c-2$, of absolute value
at most $2$, since $0<c\le1$.

Now sum \eqref{sb:comparison-function}.  The coefficient of
$y\log y$ is $(2c-2)-2c=-2$, and its constant coefficient is
\[
 (2-2c)-2c(2c-1)+4c^2=2.
\]
The error in replacing $4c^2\lfloor y\rfloor$ by $4c^2y$ is at most
$4$.  It follows that
\begin{align*}
 \left|\sum_{n\le y}f(n)
           -y\bigl((\log y)^2-2\log y+2\bigr)\right|
 \le30\sqrt y(\log y+1)+24\sqrt y+4.
\end{align*}
To compare this with $P_2(y)$, use \eqref{sb:log-square-sum} and
\[
 \log y\le\int_1^y t^{-3/4}\,dt
            =4(y^{1/4}-1)\le4y^{1/4}.
\]
In particular $(\log y)^2\le16\sqrt y$.  The total error is at most
\[
 30\sqrt y(\log y+1)+40\sqrt y+6
       \le80\sqrt y(\log y+1),
\]
as required.
\end{proof}

\subsection{Selberg's symmetry formula}
\label{sb:symmetry-subsection}

For arithmetic functions on positive integers write
\[
 (a*b)(n)=\sum_{d\mid n}a(d)b(n/d).
\]
Let $\mathbf1(n)=1$ and let $\varepsilon(1)=1$,
$\varepsilon(n)=0$ for $n>1$.  The M\"obius function $\mu$ is $0$
when a square of a prime divides $n$, and otherwise is $(-1)^r$,
where $r$ is the number of distinct prime factors of $n$.
Summing over subsets of these prime factors proves
$\mu*\mathbf1=\varepsilon$.  Also
$\tau=\mathbf1*\mathbf1$.  All convolution rearrangements below are
rearrangements of finite divisor sums.

\begin{lemma}
\label{sb:symmetry}
For $x\ge1$, put
\[
 A(x)=\sum_{n\le x}\bigl(\Lambda(n)\log n+(\Lambda*\Lambda)(n)\bigr).
\]
Then
\begin{equation}
 |A(x)-2x\log x|\le2048x.
 \label{sb:symmetry-bound}
\end{equation}
\end{lemma}

\begin{proof}
Let $L(n)=\log n$ and define $(\mathcal Da)(n)=a(n)\log n$.
Since $\log n=\log d+\log(n/d)$,
\[
 \mathcal D(a*b)=(\mathcal Da)*b+a*(\mathcal Db).
\]
The divisor identity used earlier is $L=\mathbf1*\Lambda$.  Applying
$\mathcal D$ gives the exact pointwise identity
\[
 L^2=L*\Lambda+\mathbf1*(\mathcal D\Lambda)
     =\mathbf1*(\Lambda*\Lambda+\mathcal D\Lambda).
\]
Convolution with $\mu$ therefore gives
\begin{equation}
 \mu*L^2=\Lambda*\Lambda+\mathcal D\Lambda.
 \label{sb:exact-symmetry}
\end{equation}
Similarly $\mathcal D\tau=2L*\mathbf1$.  With the finite parameter
$c$ from \eqref{sb:finite-c} and $f$ from
\eqref{sb:comparison-function},
\begin{equation}
 \mu*f=2L-2c\mathbf1+4c^2\varepsilon.
 \label{sb:comparison-convolution}
\end{equation}

For each $d\le x$, apply \eqref{sb:comparison-error} at $y=x/d$.
Since $|\mu(d)|\le1$, reversing finite sums gives
\begin{align*}
 \left|A(x)-\sum_{n\le x}(\mu*f)(n)\right|
 &\le80\sum_{d\le x}\sqrt{x/d}\bigl(\log(x/d)+1\bigr).
\end{align*}
The preceding integral bound gives
$\log z+1\le4z^{1/4}$ for $z\ge1$, while a finite rectangle
comparison gives
\[
 \sum_{d\le x}d^{-3/4}
       \le1+\int_1^x t^{-3/4}\,dt\le4x^{1/4}.
\]
Consequently the last error is at most $1280x$.

By \eqref{sb:comparison-convolution},
\[
 \sum_{n\le x}(\mu*f)(n)=2P_1(x)-2c\lfloor x\rfloor+4c^2.
\]
Its distance from $2x\log x$ is at most $10x$: indeed
\eqref{sb:log-sum}, $0<c\le1$, and $x\ge1$ bound that distance by
$6x+4\le10x$.  Thus $1290x$ already bounds
$|A(x)-2x\log x|$, and the stated constant $2048$ is valid.
The auxiliary parameter $c$ has disappeared from the conclusion.
\end{proof}

\subsection{The equation for the error term}
\label{sb:error-subsection}

\begin{lemma}
\label{sb:error-equation}
For $x\ge1$, define
\[
 E_R(x)=R(x)\log x+\sum_{n\le x}\Lambda(n)R(x/n).
\]
Then
\begin{equation}
 |E_R(x)|\le4096x.
 \label{sb:error-bound}
\end{equation}
\end{lemma}

\begin{proof}
Integrating the finite step function $\psi$ gives
\[
 \sum_{n\le x}\Lambda(n)\log n
      =\psi(x)\log x-\int_1^x\frac{\psi(t)}t\,dt.
\]
The integral is nonnegative and at most $4x$ by
\eqref{sb:chebyshev-bound}.  Also
\[
 \sum_{n\le x}(\Lambda*\Lambda)(n)
       =\sum_{n\le x}\Lambda(n)\psi(x/n).
\]
The symmetry bound therefore implies
\[
 \left|\psi(x)\log x+\sum_{n\le x}\Lambda(n)\psi(x/n)
                -2x\log x\right|\le2052x.
\]
Substitute $\psi(y)=R(y)+y$.  The difference between the expression
inside this absolute value and $E_R(x)$ is
\[
 x\left(\sum_{n\le x}\frac{\Lambda(n)}n-\log x\right),
\]
whose absolute value is at most $6x$.  Hence
$|E_R(x)|\le2058x\le4096x$.
\end{proof}

\subsection{A second substitution and finite summation by parts}
\label{sb:weighted-sum-subsection}

\begin{lemma}
\label{sb:weighted-sum}
For $x\ge4$,
\begin{equation}
 |R(x)|(\log x)^2
 \le2\sum_{n\le x}\log n\,|R(x/n)|+2^{17}x\log x.
 \label{sb:weighted-sum-bound}
\end{equation}
\end{lemma}

\begin{proof}
Multiply the defining equation for $E_R(x)$ by $\log x$, split
$\log x=\log n+\log(x/n)$ in its sum, and use the same equation at
$x/n$.  This gives the exact identity
\begin{align*}
 R(x)(\log x)^2={}&
   -\sum_{n\le x}\Lambda(n)\log n\,R(x/n)\\
 &+\sum_{ab\le x}\Lambda(a)\Lambda(b)R(x/(ab))\\
 &+E_R(x)\log x-\sum_{n\le x}\Lambda(n)E_R(x/n).
\end{align*}
The last line has absolute value at most
\[
 4096x\left(\log x+\sum_{n\le x}\frac{\Lambda(n)}n\right)
 \le4096x(2\log x+6)
 \le32768x\log x,
\]
because $\log x\ge\log4\ge1$.  The arithmetic function
\[
 a(n)=\Lambda(n)\log n+(\Lambda*\Lambda)(n)
\]
is nonnegative.  Taking absolute values gives
\begin{equation}
 |R(x)|(\log x)^2
    \le\sum_{n\le x}a(n)|R(x/n)|+32768x\log x.
 \label{sb:first-weighted-bound}
\end{equation}

Put $N=\lfloor x\rfloor$, $g_n=|R(x/n)|$ for $1\le n\le N+1$,
and
\[
 \Delta_n=A(n)-2P_1(n).
\]
The symmetry estimate and \eqref{sb:log-sum} give
\[
 |\Delta_n|\le2048n+4n\le4096n.
\]
Since $x/(N+1)<1$, we have $g_{N+1}=x/(N+1)$.  The exact finite
summation-by-parts identity is
\begin{equation}
 \sum_{n=1}^N(a(n)-2\log n)g_n
   =\Delta_Ng_{N+1}
       +\sum_{n=1}^N\Delta_n(g_n-g_{n+1}).
 \label{sb:finite-abel}
\end{equation}
Monotonicity of $\psi$ gives
\[
 |g_n-g_{n+1}|
 \le\psi(x/n)-\psi(x/(n+1))+\frac{x}{n(n+1)}.
\]
In the first term of the resulting weighted sum, each $\Lambda(m)$
is counted with weight $\lfloor x/m\rfloor\le x/m$.  Therefore
\begin{align*}
 \sum_{n=1}^N n|g_n-g_{n+1}|
 &\le x\sum_{m\le x}\frac{\Lambda(m)}m
                   +x\sum_{n=1}^N\frac1{n+1}\\
 &\le x(\log x+6)+x\log(N+1)
 \le x(2\log x+7).
\end{align*}
Here $N+1\le2x$ and $\log2\le1$.  The endpoint term in
\eqref{sb:finite-abel} is at most $4096x$.  Consequently
\[
 \left|\sum_{n\le x}(a(n)-2\log n)|R(x/n)|\right|
 \le4096x(2\log x+8)\le40960x\log x.
\]
Combine this with \eqref{sb:first-weighted-bound}.  The two error
constants add to $73728$, which is less than $2^{17}$.
\end{proof}

\subsection{Replacing the weighted sum by a bounded integral}
\label{sb:sum-integral-subsection}

\begin{lemma}
\label{sb:sum-integral}
For $x\ge1$,
\begin{equation}
 \left|\sum_{n\le x}\log n\,|R(x/n)|
            -\int_1^x\log t\,|R(x/t)|\,dt\right|\le20x.
 \label{sb:sum-integral-bound}
\end{equation}
\end{lemma}

\begin{proof}
Set $N=\lfloor x\rfloor$ and, for $t\ge1$, put
$g(t)=|R(x/t)|$ and $h(t)=\log t\,g(t)$.  The bound
$g(t)\le4x/t$ also holds when $t>x$, because then $R(x/t)=-x/t$.
For $n\le t\le n+1$, monotonicity of $\psi$ gives
\begin{align*}
 |h(t)-h(n)|
 \le{}&\log n\left(\psi(x/n)-\psi(x/t)+\frac xn-\frac xt\right)\\
      &+(\log t-\log n)g(t).
\end{align*}
Integrate over each of the $N$ unit intervals and sum.

For the contribution involving $\psi$, each $\Lambda(m)$ can occur
only in the one unit interval containing $x/m$.  Its contribution is
at most $\Lambda(m)\log(x/m)$, because that interval has length at
most $1$ and $n\le x/m$.  The total is thus at most
\[
 \sum_{m\le x}\Lambda(m)\log(x/m)
      =\int_1^x\frac{\psi(u)}u\,du\le4x.
\]
The contribution from $x/n-x/t$ is at most
\[
 x\sum_{n=1}^N\frac{\log n}{n(n+1)}\le x.
\]
To verify the last inequality using only finite sums, telescope:
\begin{align*}
 \sum_{n=1}^N\log n\left(\frac1n-\frac1{n+1}\right)
 &=\sum_{n=2}^N\frac{\log n-\log(n-1)}n
                         -\frac{\log N}{N+1}\\
 &\le\sum_{n=2}^N\frac1{n(n-1)}\le1.
\end{align*}
Finally, $\log t-\log n\le(t-n)/n$ and $g(t)\le4x/t$ bound the
last contribution by
\[
 4x\sum_{n=1}^N\frac1{n^2}\le8x.
\]
Indeed $1/n^2\le1/(n(n-1))$ for $n\ge2$, so the last sum is at
most $2$.  We have proved
\[
 \left|\sum_{n=1}^Nh(n)-\int_1^{N+1}h(t)\,dt\right|\le13x.
\]
On the interval $[x,N+1]$, $h(t)=x\log t/t\le x$, and the length
of the interval is at most $1$.  Removing that final integral gives
the stronger bound $14x$, hence the displayed bound $20x$.
\end{proof}

\begin{theorem}
\label{sb:contraction-inputs}
Write $r(t)=e^{-t}R(e^t)$ for $t\ge0$.  Then
\begin{align}
 |r(t)|&\le4,
 \label{sb:r-bound}\\
 \left|\int_a^b r(v)\,dv\right|&\le20
          &&(0\le a\le b),
 \label{sb:r-signed}\\
 |r(t)|&\le\frac2{t^2}\int_0^t(t-v)|r(v)|\,dv+\frac{2^{18}}t
          &&(t\ge2).
 \label{sb:r-weighted}
\end{align}
For $0\le u\le v$ one also has
\begin{equation}
 r(v)\ge r(u)-4(v-u).
 \label{sb:r-downward}
\end{equation}
Thus, in the contraction lemma, one may use
$B=4$, $C=20$, $S=2^{18}$, and the harmless larger downward constant
$L=5$.
\end{theorem}

\begin{proof}
The first assertion is \eqref{sb:chebyshev-bound}; the second is
\eqref{sb:signed-integral} after the finite change of variables
$x=e^v$.  Lemmas~\ref{sb:weighted-sum} and \ref{sb:sum-integral} give,
for $x\ge4$,
\[
 |R(x)|(\log x)^2
 \le2\int_1^x\log u\,|R(x/u)|\,du
                    +2^{17}x\log x+40x.
\]
Since $\log x\ge1$, the last two terms are at most
$2^{18}x\log x$.  Put $t=\log x$ and $v=\log(x/u)$.  Direct
substitution yields
\[
 \frac1x\int_1^x\log u\,|R(x/u)|\,du
       =\int_0^t(t-v)|r(v)|\,dv.
\]
Divide by $t^2$ to obtain \eqref{sb:r-weighted}; $t\ge2$ ensures
$e^t\ge4$.

Finally, monotonicity of $\psi$ implies
\[
 r(v)+1=\frac{\psi(e^v)}{e^v}
       \ge e^{-(v-u)}\frac{\psi(e^u)}{e^u}
       =e^{-(v-u)}(r(u)+1).
\]
Because $0\le r(u)+1\le4$ and $1-e^{-h}\le h$ for $h\ge0$,
this gives \eqref{sb:r-downward}.
\end{proof}

\section{A finite contraction with an explicit modulus}
\label{con:section}

This section uses the estimates proved above, with deliberately generous
constants.  Put
\[
 r(t)=e^{-t}\bigl(\psi(e^t)-e^t\bigr),\qquad S=2^{18}.
\]
The hypotheses needed here are
\begin{align}
 |r(t)|&\leq 4 &&(t\geq0),\label{con:global}\\
 \left|\int_u^v r(w)\,dw\right|&\leq20
       &&(0\leq u\leq v),\label{con:signed}\\
 |r(t)|&\leq \frac{2}{t^2}\int_0^t(t-u)|r(u)|\,du+\frac{S}{t}
       &&(t\geq2).\label{con:weighted}
\end{align}
There is also a useful one-sided regularity estimate.  Monotonicity of
$\psi$ gives, for $v\geq u$,
\[
 r(v)+1\geq e^{u-v}(r(u)+1).
\]
Since $0\leq\psi(e^u)e^{-u}\leq4$ and $1-e^{-h}\leq h$, it follows that
\begin{equation}
 r(v)\geq r(u)-5(v-u).\label{con:decline}
\end{equation}
The extra unit in the constant is harmless.  In particular, no two-sided
continuity assumption is being made: $r$ has upward jumps.
At rational spatial points $x\le y$, the same calculation is simply
$\psi(y)/y\ge\psi(x)/y$ together with $1-x/y\le\log(y/x)$;
no exponential evaluation is needed in its arithmetic implementation.

\subsection{A deficit detected by finitely many rational tests}

Fix an integer $q\geq1$, and define
\begin{equation}
 \varepsilon=\frac1q,\quad D=2^{20}q^3,\quad
 \delta=\frac1D,\quad \overline H=1024q,\quad
 Q=1+\frac1{320q},\quad m=327680q^2,\quad
 \sigma=\frac1{64q}.
 \label{con:parameters}
\end{equation}
Both $m$ and $m/2$ are integers.  Write $\lambda=\log Q$ and
$H=m\lambda$.  The elementary inequalities
\[
 \frac{z}{1+z}\leq\log(1+z)\leq z\qquad(z\geq0)
\]
imply
\begin{equation}
 \frac1{640q}\leq\lambda\leq\frac1{320q},
 \qquad 512q\leq H\leq1024q=\overline H.
 \label{con:block-length}
\end{equation}

\begin{lemma}[Finite block deficit]\label{con:deficit}
Suppose $\varepsilon\leq a\leq4$ and $|r(t)|\leq a$ throughout a
block $[s,s+H]$, where $e^s$ is rational.  Then
\begin{equation}
 \int_s^{s+H}\bigl(a-|r(t)|\bigr)\,dt
       \geq\frac1{64q^2}.
 \label{con:deficit-bound}
\end{equation}
The proof needs only finitely many comparisons of rational numbers.
\end{lemma}

\begin{proof}
The nodes
\[
 x_i=e^sQ^i,\qquad t_i=\log x_i=s+i\lambda,
       \qquad 0\leq i\leq m,
\]
have rational $x_i$.  Compute rational numbers $b_i$ satisfying
\begin{equation}
 |b_i-r(t_i)|\leq\sigma.\label{con:sample-error}
\end{equation}
Such numbers are obtained by a bounded prime-power sum and rational
logarithm enclosures.  Indeed, approximating each nonzero $\Lambda(n)$
with absolute error at most $\sigma$ produces an error at most
$\sigma\lfloor x_i\rfloor/x_i\leq\sigma$ after division by $x_i$.
All tests below concern the rational numbers $b_i$ and $a/4$.

Set $d(t)=a-|r(t)|$.  First suppose that
$b_i<a/4$ for all $0\leq i\leq m/2$.
For $t\in[t_i,t_{i+1}]$ in the first half of the block,
\eqref{con:decline}, applied from $t$ to $t_{i+1}$, gives
\[
 r(t)\leq r(t_{i+1})+5\lambda
       \leq a/4+\sigma+5\lambda
       \leq 9a/32<a/2.
\]
For $-a\leq z\leq a/2$ one has
$a-|z|\geq(a+z)/3$.  Therefore \eqref{con:signed} and
\eqref{con:block-length} imply
\[
 \int_s^{s+H/2}d(t)\,dt
 \geq\frac{aH/2-20}{3}
 \geq\frac{256-20}{3}
 \geq\frac1{64q^2}.
\]
This proves the claim in the first case.

Second, suppose that $b_i>-a/4$ for all $m/2\leq i\leq m$.
For each interval in the second half, apply \eqref{con:decline}
from its left endpoint to obtain
$r(t)\geq-a/4-\sigma-5\lambda\geq-9a/32>-a/2$.
The inequality $a-|z|\geq(a-z)/3$ for $-a/2\leq z\leq a$
gives the same lower bound for the integral of $d$ on this half.

If neither case applies, a bounded search finds an index $i_0$ in the
first half with $b_{i_0}\geq a/4$ and an index in the second half with
$b_i\leq-a/4$.  Let $v$ be the first index after $i_0$ with
$b_v\leq-a/4$, and let $u<v$ be the last index at or after $i_0$
with $b_u\geq a/4$.  These are minima and maxima of finite decidable
sets.  For $u<i<v$,
\[
 -a/4<b_i<a/4.
\]
For each $i$ from $u$ through $v-1$, the left sample is at least
$-a/4$ and the right sample is at most $a/4$.  Thus the two applications
of \eqref{con:decline} used in the preceding cases yield
\[
 |r(t)|\leq a/4+\sigma+5\lambda\leq9a/32<a/2
       \qquad(t_u\leq t\leq t_v).
\]
At the outer samples,
\[
 r(t_u)-r(t_v)\geq a/2-2\sigma\geq15a/32.
\]
Another application of \eqref{con:decline} gives
$t_v-t_u\geq15a/(32\cdot5)$.  Consequently
\[
 \int_{t_u}^{t_v}d(t)\,dt
 \geq\frac a2\frac{15a}{32\cdot5}
 =\frac{3a^2}{64}\geq\frac1{64q^2}.
\]
All three cases prove \eqref{con:deficit-bound}.  No sign decision for
an arbitrary real number and no intermediate-value principle occurs.
\end{proof}

\subsection{Using the deficit in the weighted inequality}

\begin{lemma}[One finite contraction]\label{con:one-step}
Let $s,t\geq0$ with $e^s$ and $e^t$ rational, and suppose
$\varepsilon\leq a\leq4$ and $|r(u)|\leq a$ for $s\leq u\leq t$.
If
\begin{equation}
 t\geq8(s+\overline H),\qquad
 t\geq D(8s+S),\qquad t\geq2,
 \label{con:step-threshold}
\end{equation}
then
\begin{equation}
 |r(t)|\leq a-3\delta.\label{con:step-margin}
\end{equation}
\end{lemma}

\begin{proof}
Place consecutive full blocks of length $H$ starting at $s$ and ending
no later than $t/2$.  Their number $b$ satisfies
\[
 b\geq \frac{t/2-s}{H}-1
   \geq\frac{3t}{8H}-1
   \geq\frac{t}{4H}
   \geq\frac{t}{4\overline H}.
\]
Here $s\leq t/8$ and $t\geq8\overline H\geq8H$ were used.
There is no need to decide a floor of a real number: when $x=e^t$ is
rational, the permitted blocks are found by the rational test
\[
 \bigl(e^s Q^{m i}\bigr)^2\leq x.
\]
The number of candidates can be bounded by $\lceil x\rceil+1$:
Bernoulli's finite inequality gives
$Q^{mi}\geq1+mi/(320q)=1+1024qi>i$.

On all these blocks the weight $t-u$ is at least $t/2$.
By Lemma~\ref{con:deficit},
\begin{align*}
 \int_0^t(t-u)|r(u)|\,du
 &\leq\frac{at^2}{2}+4st
      -b\frac{t}{2}\frac1{64q^2}\\
 &\leq\frac{at^2}{2}+4st
      -\frac{t^2}{8\overline H}\frac1{64q^2}.
\end{align*}
The term $4st$ bounds the possible excess above the baseline $a$
on $[0,s]$, using \eqref{con:global}.
Substitution in \eqref{con:weighted} gives
\[
 |r(t)|\leq a
    -\frac1{4\overline H\cdot64q^2}
    +\frac{8s+S}{t}
 =a-4\delta+\frac{8s+S}{t}
 \leq a-3\delta.
\]
This proves the stated, deliberately stronger, decrement.
\end{proof}

Every part of this argument refers only to arguments no larger than $x$.
This finite dependence is essential in the induction below.  The surplus
$2\delta$ in \eqref{con:step-margin}, beyond the decrement $\delta$
that will actually be used, also provides room for certified numerical
enclosures.  The rational-arithmetic implementation assigns an explicit
error budget to the compound expressions; the coarse sample precision
$\sigma$ in Lemma~\ref{con:deficit} is not the precision used for those
compound expressions.

\subsection{Integer thresholds and the modulus}

Set
\begin{equation}
 J=4D-D/q,\qquad a_j=4-j/D\quad(0\leq j\leq J),
 \label{con:levels}
\end{equation}
so that $J$ is an integer and $a_J=1/q$.  Define, by primitive recursion,
\begin{align}
 T_0&=2,\notag\\
 T_{j+1}&=16(T_j+\overline H)+D(16T_j+S)+1,\label{con:recurrence}\\
 P_j&=2^{2T_j}.\notag
\end{align}
These are integers.  Since $1/2\leq\log2\leq1$,
\begin{equation}
 T_j\leq\log P_j\leq2T_j.\label{con:dyadic-threshold}
\end{equation}
Starting at the rational point $P_j$, the hypotheses of
Lemma~\ref{con:one-step} are therefore met at every rational
$x\geq P_{j+1}$.  Indeed, putting $s=\log P_j$ and $t=\log x$,
\eqref{con:recurrence} bounds both $8(s+\overline H)$ and
$D(8s+S)$ from above, with a spare unit.
We retain only a decrement $\delta$ at each step.  After $J$ steps,
\begin{equation}
 |\psi(x)-x|\leq x/q\qquad(x\geq P_J).
 \label{con:psi-at-recursive-threshold}
\end{equation}
The arithmetic justification of this finite induction is given in the
next subsection.

The recurrence admits a simple closed upper bound.  Put $F=2^{20}D$.
From \eqref{con:recurrence} and $D=2^{20}q^3$ one checks
\[
 T_{j+1}+1\leq F(T_j+1),\qquad
 T_J+1\leq3F^J.
\]
For example, the coefficient of $T_j$ in the recurrence is
$16+16D<F$, and its constant term plus one is
$2^{14}q+2^{18}D+2<F$.
The elementary inequality $2^n\geq n^2$ for $n\geq4$ gives
$D\leq2^{D/16}$, because $D\geq2^{20}$.  Hence
\[
 F=2^{20}D\leq2^{D/8},\qquad
 3F^J\leq3\,2^{D^2/2}\leq2^{D^2},
\]
where $J\leq4D$ was used.  Thus $T_J\leq2^{D^2}$ and
\begin{equation}
 P_J\leq N_\psi(q):=2^{\,2^{D^2+1}},
       \qquad D=2^{20}q^3.
 \label{con:closed-modulus}
\end{equation}
In particular, the modulus is elementary, and therefore primitive
recursive.  No claim of numerical usefulness is intended.

\subsection{The induction predicate is a bounded finite check}
\label{con:bounded-induction}

Induction on the assertion ``for every $x\geq P_j$'' would not be an
instance of quantifier-free induction.  Instead, fix a final integer
ceiling $X$, as well as $q$, before starting the induction.  Form
rational lower and upper enclosures
$\underline\psi(n)\leq\psi(n)\leq\overline\psi(n)$
for $0\leq n\leq X$.  The stage predicate is the Boolean conjunction,
over the finite set of integers $P_j\leq n\leq X$, of
\begin{equation}
 \overline\psi(n)\leq(1+a_j)n,
 \qquad
 \underline\psi(n)\geq(1-a_j)(n+1).
 \label{con:endpoint-predicate}
\end{equation}
Denominators are cleared, so this is a primitive-recursive Boolean
function of its integer arguments.  Its equality to $1$ is a
quantifier-free formula.

The two checks enclose the value at the left of the cell $[n,n+1)$
and its limit at the right. Between consecutive integers, $\psi(x)$ is
constant and $\psi(x)/x-1$ decreases. When $a_j>1$ the lower bound
is also immediate from positivity. Therefore
\eqref{con:endpoint-predicate} controls every rational argument needed
in the next contraction whose value is less than $X+1$.
All mesh coordinates in Lemma~\ref{con:deficit} are rational, and their
indices and denominators have explicit primitive-recursive bounds.
An empty range in \eqref{con:endpoint-predicate} is checked as true.

The base envelope $a_0=4$ has slack: $0\leq\psi(x)\leq4x$ implies
$-1\leq r(t)\leq3$.  For the step, the analytic bound improves $a_j$
by $3\delta$, whereas the new check asks for improvement only by
$\delta$. In particular, each normalized logarithm enclosure may be
required to have error at most $\delta/8$. The lower end-of-cell
condition is tested at $z=n+1-\delta/64$, where $\psi(z)=\psi(n)$.
This point is at least $n$ and less than $X+1$, so both its threshold
and the finite ceiling are preserved. Section~\ref{sec:arithmetic}
gives the explicit algebra recovering the end-of-cell condition.
The remaining compound-expression errors are assigned a total budget
at most $\delta/4$ by the finite arithmetic compilation described
elsewhere in this report.  The surplus in
\eqref{con:step-margin} exceeds all these allowances.

Thus the stage implication is a finite arithmetic calculation with
parameters $q$ and $X$, and quantifier-free induction applies to its
Boolean check.  Only after the $J$ stages have been completed is the
arbitrary ceiling $X$ generalized.  This is the reason the proof does
not use an induction hypothesis with an unbounded universal quantifier.


\section{From the Chebyshev estimate to prime counting}\label{sec:pi}

Let $\vartheta(x)=\sum_{p\le x}\log p$. The preceding sections prove, for
each integer $q\ge1$,
\begin{equation}\label{eq:psi-modulus}
 \abs{\frac{\psi(x)}x-1}\le\frac1q
 \quad\text{when}\quad
 x\ge N_\psi(q):=2^{\,2^{(2^{20}q^3)^2+1}}.
\end{equation}
We now prove the conversion with constants; it uses only finite sums and
the global bound $\vartheta(u)\le\psi(u)\le4u$.

\begin{lemma}\label{lem:pi-conversion}
For $x\ge4$ satisfying \eqref{eq:psi-modulus},
\begin{equation}\label{eq:pi-error}
 \abs{\frac{\pi(x)\log x}{x}-1}
 \le \frac1q+\frac{24\log x}{\sqrt{x}}+\frac{16}{\log x}.
\end{equation}
\end{lemma}
\begin{proof}
Prime powers of exponent at least two contribute
\[
 0\le\psi(x)-\vartheta(x)
 =\sum_{a=2}^{\floor{\log x/\log2}}\vartheta(x^{1/a})
 \le4\sqrt{x}\frac{\log x}{\log2}
 \le8\sqrt{x}\log x,
\]
since $\log2\ge1/2$. Summation by parts gives the exact finite identity
\[
 \pi(x)=\frac{\vartheta(x)}{\log x}
       +\int_2^x\frac{\vartheta(u)}{u(\log u)^2}\,du.
\]
On $[2,\sqrt{x}]$ the integrand is at most
$4/(\log2)^2\le16$. On $[\sqrt{x},x]$ it is at most
$16/(\log x)^2$. Thus
\[
 0\le\pi(x)\log x-\vartheta(x)
 \le16\sqrt{x}\log x+\frac{16x}{\log x}.
\]
Adding the two displayed bounds to the error for $\psi$ proves the claim.
The integral here can also be written as a finite telescoping sum, since
the antiderivative of $1/(u\log^2u)$ is $-1/\log u$.
\end{proof}

The already proved inequality $\log x\le4x^{1/4}$ gives
$(\log x)^2\le16\sqrt{x}$. Thus, with $u=\log x$, the last two terms
of \eqref{eq:pi-error} are at most $400/u$.
With $q=8k$ and $u\ge4096k$, this gives
\begin{equation}\label{eq:pi-quarter}
 \abs{\frac{\pi(x)\log x}{x}-1}
 \le\frac1{8k}+\frac{400}{4096k}
 =\frac{57}{256k}<\frac1{4k}.
\end{equation}

For $n\ge N(k)$ from \eqref{eq:modulus}, we have
$n\ge N_\psi(8k)$, $\log n\ge4096k$, and $n\ge4k$.
For example,
$\log N(k)\ge 2^{D^2+1}\ge4096k$ follows from
$\log2\ge1/2$ and $D=2^{20}(8k)^3$.
Since $\pi(n)\le n$, replacing $\log n$ by $L(n)$ costs at most $1/n$:
\begin{equation}\label{eq:rational-half}
 \abs{\frac{\pi(n)L(n)}n-1}
 <\frac1{4k}+\frac1n\le\frac1{2k}.
\end{equation}
Clearing the positive denominator $nF(n)$ proves the stronger inequality
\[
 2k\,\dist(\pi(n)A(n),nF(n))\le nF(n)
 \qquad(n\ge N(k)).
\]
In particular, it proves the announced matrix with coefficient $k$.
The factor-two surplus is useful for the finite numerical enclosures below.

\section{Why this is a proof in iPRA}\label{sec:arithmetic}

This section supplies the logical bookkeeping. Merely knowing that a modulus
is primitive recursive would not suffice. We must show that its correctness
is established with the induction and logic actually available in $\ipra$.

\subsection{Bounded tests and finite arithmetic}

Represent integers by a sign and a natural magnitude, and rationals by
integer numerator and positive natural denominator. Addition, multiplication,
comparison and exact division are primitive recursive by their elementary
bounded algorithms. Canonical reduction of a fraction is optional. Factorials,
finite sums and products, bounded division, powers, prime factorization by
successive bounded trial division, and the functions $\mu,\tau,\pi$ are all
primitive recursive. Their identities used above are proved by induction on
bounded lists of factors, sums or products.

For any fixed primitive-recursive Boolean $b(i,\vec z)$, define
\[
 \operatorname{All}(0,\vec z)=b(0,\vec z),\qquad
 \operatorname{All}(u+1,\vec z)
 =\operatorname{All}(u,\vec z)\mathbin{\&}b(u+1,\vec z).
\]
Quantifier-free induction proves that $\operatorname{All}(u,\vec z)=1$
is equivalent to the bounded assertion that $b(i,\vec z)=1$ for every
$i\le u$. In the extraction direction, keep the queried index $i$ as a
free parameter and induct on the quantifier-free implication
$i\le u\to(\operatorname{All}(u,\vec z)=1\to b(i,\vec z)=1)$;
generalize $i$ only afterward. Bounded existential tests use the
corresponding Boolean disjunction and bounded witness search.
This applies also to nested bounded quantifiers after their inner tests have
been compiled. Every induction over a table in this proof is induction on
such a Boolean equation. It is not induction on a formula with an unbounded
quantifier hidden inside it.

\subsection{A common rational table for the entire finite argument}

Fix integers $k\ge1$ and $X\ge2$, and take $q=8k$.
For a precision parameter $M$ define, for every prime $p\le X+1$,
\[
 \ell_p=H_{Mp}-H_M,\qquad u_p=\ell_p+\frac1M.
\]
For integers $n\le X+1$, let $c(n)$ be the number of prime powers at most
$n$, counted once for each pair $(p,a)$ with $p$ prime and $p^a\le n$.
Set
\[
 L_\psi(n)=\sum_{p^a\le n}\ell_p,\qquad
 U_\psi(n)=L_\psi(n)+\frac{c(n)}M.
\]
There are at most $n$ such prime powers, so
\begin{equation}\label{eq:psi-table}
 L_\psi(n)\le\psi(n)\le U_\psi(n),\qquad
 U_\psi(n)-L_\psi(n)\le n/M.
\end{equation}
All sums in this display are bounded; the displayed real $\psi$ is only
the explanatory reading of its rational enclosures.

With $D,\delta=1/D,a_j,P_j,J$ as in the contraction section, the stage
predicate is the single Boolean conjunction, over $P_j\le n\le X$, of
\begin{equation}\label{eq:stage-table}
 U_\psi(n)\le(1+a_j)n,
 \qquad L_\psi(n)\ge(1-a_j)(n+1).
\end{equation}
Clear the denominators to make both tests comparisons of integers. On each
cell $n\le y<n+1$, $\psi(y)=\psi(n)$, and $\psi(n)/y$ is decreasing.
Thus \eqref{eq:stage-table} controls the entire cell, including its two
one-sided endpoint values. When $a_j>1$, the lower estimate is automatic
from positivity; otherwise it follows by dividing by $n+1$.
No quantification over arbitrary real numbers or rational denominators is
needed in the induction predicate.

At stage zero $a_0=4$. The global estimate $\psi(n)\le4n$ and $M\ge1$
give $U_\psi(n)\le5n$, and the lower inequality is immediate. Hence the
base of induction holds with room to spare.

\subsection{How the finite stage implication preserves the table}

The contraction calculation gives the exact improvement $a_j-3\delta$,
while the next predicate claims only $a_{j+1}=a_j-\delta$.
Its proof uses only arguments in the cells already controlled by the table.
For the upper endpoint take $x=n$. For the lower endpoint take the rational
number
\[
 z=n+1-\tau,\qquad \tau=\delta/64.
\]
Then $\psi(z)=\psi(n)$ and $z<X+1$ even when $n=X$.
There is no appeal to a limit at an integer discontinuity.

Suppose first that the exact contraction inequalities are used. If
$1/M\le\delta$, then the upper estimate gives
\[
 U_\psi(n)\le(1+a_j-3\delta)n+n/M
 \le(1+a_j-2\delta)n.
\]
The lower estimate gives
\[
 L_\psi(n)\ge(1-a_j+3\delta)(n+1-\tau)-n/M.
\]
After subtracting $(1-a_j+\delta)(n+1)$, its surplus is at least
\[
 2\delta(n+1)-2\tau-\delta n>0,
\]
because the positive part of $1-a_j+3\delta$ is at most $2$.
Thus both comparisons in the next table follow.
The same calculation remains valid if the finite numerical reconstruction
costs at most $\delta/4$ in the normalized contraction estimate and if
$1/M\le\delta/4$; for example the upper error then totals at most
$\delta/2$, well below the $2\delta$ reserve. This is the tolerance used next.

\subsection{Explicit finite precision and quadrature budget}\label{sec:budget}

Here are concrete bounds showing that the tolerance can be met by primitive
recursion. They also specify a straightforward route for the Lean arithmetic
library. A sharper numerical implementation is unnecessary.

Let $S=2^{18}$, $J=4D-D/q$, and put
\begin{equation}\label{eq:precision}
 W=X+q+D+S+J+16,\qquad B=2^{W^4},\qquad
 Q_{\rm quad}=2^{B^3},\qquad M=2^{Q_{\rm quad}^3}.
\end{equation}
Use the harmonic intervals of width at most $1/M$ for every logarithm and
the rational spatial mesh of width at most $1/Q_{\rm quad}$ for every finite integral.
Include all integer breakpoints and rational endpoints in that mesh. The
geometric mesh used for the deficit argument is a different, exact rational
mesh; its sample values use the same harmonic intervals.

We justify a generous uniform error bound rather than suppressing the
numerical reconstruction behind the word ``computable.'' On $[1,Y]$, with
$Y=X+1$, the total variation of $R$ is at most $5Y$ and $|R|\le4Y$.
For a decreasing nonnegative function $w$, the elementary product estimate
\[
 \operatorname{Var}(w|R|)
 \le \sup w\,\operatorname{Var}(R)
       +\sup|R|\,\operatorname{Var}(w)
\]
follows from the triangle inequality on each finite partition. Here a
variation bound means a bound on every finite partition sum; no supremum
object is introduced. In particular,
for $w(u)=\log(Y/u)/u^2$,
\begin{equation}\label{eq:var}
 \operatorname{Var}_{[1,Y]}(w|R|)\le9Y^2.
\end{equation}
A rectangle sum on mesh $1/Q_{\rm quad}$ therefore differs from its finite integral
by at most $10Y^2/Q_{\rm quad}$, including a final partial cell. This follows by summing
the cell oscillations times their widths. It does not assume an ineffective
Riemann-integrability theorem.

All the other integral expressions in the proof are simpler: sums of a
step function times $1/u$, $1/u^2$, $\log u$, $1/(u\log^2u)$, or a
polynomial in $\log u$ of degree at most two. The interval with
$1/\log u$ begins at $2$. Split at integer breakpoints, use
$\log2\ge1/2$, and apply the same finite product estimate. Their variation,
normalized coefficient sums, and the reciprocals of their nonzero
denominators are each bounded by $W^{12}<B$. For the signed integrals one
can instead use their exact finite Abel sums, eliminating quadrature there.
The substitution $t=\log y$ in the contraction uses only rational spatial
endpoints: $r(t)\,dt=R(y)\,dy/y^2$. Likewise, its weighted integral is
\[
 \int_0^{\log x}(\log x-t)|r(t)|\,dt
 =\int_1^x\frac{\log(x/y)|R(y)|}{y^2}\,dy.
\]
Thus it requires neither the computation of $\floor{e^t}$ nor decisions
at a numerically uncertain discontinuity.

For clarity, the coarse operation and size audit behind $W^{12}$ is as follows.
One divisor or prime-power table uses at most $(Y+1)^2$ entries; a double
convolution uses at most $(Y+1)^3$ elementary table accesses. A geometric
block uses $327680q^2$ samples and there are at most $Y+1$ relevant blocks.
There are at most $X+1$ cell checks at each of at most $J+1$ stages.
The remaining finite algebra has fixed degree. Logs have magnitude at most
$2(Y+1)^2$, including factorial and binomial intermediates and log ratios;
division by a block
step costs at most $640q$. These bounds are all dominated by $W^{12}$.
Thresholds exceeding $X+1$ require no analytic evaluation: their table
conditions are empty, decided by integer comparison. Large numerators and
denominators of exact rational coordinates affect running time, but not
the magnitudes or the enclosure error estimates.

Treat exact rational arithmetic as exact, and each logarithm or finite
integral evaluation as an interval operation with its proved error bound.
Across an addition, multiplication, absolute value or reciprocal whose
arguments have magnitude at most $B$ and whose required divisor has
absolute value at least $1/B$, an accumulated error is multiplied by at
most $8B^3$, provided the input error is at most $1/(2B)$.
The reciprocal estimate is
\[
 \abs{\frac1u-\frac1v}
 \le 2B^2|u-v|\qquad
 (|u|\ge1/B,\ |u-v|\le1/(2B)).
\]
Finite summation uses the sum of the coefficient bounds. A quadrature
operation has coefficient sum at most its interval length; consequently its
error bound does not grow with the number of mesh cells. It may be
treated as one proved bounded-sum operation in this audit.
There are fewer than $B$ such error-bearing operations along the complete
finite argument. The exact internal summands of a harmonic or quadrature
macro are accounted for by that macro's separately proved bound.
The total propagated error is therefore bounded by
\begin{equation}\label{eq:budget-bound}
 2B(8B^3)^B\left(\frac1{Q_{\rm quad}}+\frac1M\right)
 \le\frac{4B(8B^3)^B}{Q_{\rm quad}}
 <\min\left\{\frac1{64D},\frac1{2B}\right\}.
\end{equation}
Indeed $B\ge16$, $D<W<B$, and the base-two logarithm of the numerator is
less than $6B^2$, whereas $\log_2 Q_{\rm quad}=B^3$. The choice of $B$ is enormously
larger than needed. It also ensures the reciprocal side condition and the
sample tolerance $1/(64q)$ used in the deficit lemma.

This argument is a finite error propagation calculation: its invariant is
the rational bound on each entry of a finite computation table. Its induction
is an equation or inequality between primitive-recursive terms. The log and
quadrature macros themselves are bounded sums with the rectangle estimates
proved above. The long harmonic sums are exact rational computations; we
do not add a rounding error for each of their summands. The quadrature mesh
is chosen first; the precision of logarithm evaluations is chosen afterward.
This distinction prevents a circular precision estimate.

In particular, every occurrence of analytic notation in the stage implication
can be replaced by rational intervals with total loss less than
$\delta/64$. This is smaller than the $\delta/4$ allowed in the preceding
endpoint calculation. The decisions in the deficit lemma are made on the
rational sample approximations, with separated thresholds, so they remain
valid at this precision. The displayed Selberg constants also have ample
reserve from their preceding sharper inequalities. Thus the implication
between consecutive Boolean stage tables is a derivation in $\ipra$,
using only the stated defining equations and quantifier-free induction.

\begin{lemma}[Translation of the finite calculation schema]\label{lem:translation}
The finite calculation rules used here have uniform $\ipra$ derivations.
In particular, replacing their displayed logarithms and integrals by the
specified rational intervals produces a derivation of the implication
between consecutive stage tables, not merely an algorithm that computes
the two tables.
\end{lemma}
\begin{proof}
For exact arithmetic, prove the ring and order identities by induction on
the defining primitive recursions; represent signed and rational identities
by natural-number identities after clearing positive denominators.
The identity for a finite sum, and an inequality obtained by summing its
termwise inequalities, are proved by induction on the number of summands.
For a bounded search, induction on the scanned prefix proves the invariant
that it contains the first, or last, successful test, if one has occurred.
These are Boolean invariants and require no unbounded least-number rule.

The rectangle, refinement, finite Abel and error-propagation rules are also
finite-sum rules. Their proofs have been given by termwise inequalities and
telescoping; apply the same bounded-sum induction to those proofs. For a
loop, conjoin the row invariants through its prescribed finite bound,
using $\operatorname{All}$, and apply quantifier-free induction to that
conjunction. Nested loops are handled from the inside outward. All bounds
are displayed primitive-recursive terms or compositions of them. Ordinary
intuitionistic implication and universal introduction retain their free
integer parameters throughout.

Finally, the three cases in the deficit proof are exhaustive Boolean
alternatives. Their exhaustive character is proved by the bounded-search
invariant, not by excluded middle for a real predicate. Each branch was
proved by the above termwise rules, with positive margins larger than the
enclosure losses. The threshold and endpoint calculations are exact
rational inequalities. Composing these finitely many rule derivations
therefore gives the required stage implication in $\ipra$.
This translation uses the proofs of the rows and loop invariants; it does
not infer the correctness of an arbitrary primitive-recursive program from
its computability or from tests in the standard model.
\end{proof}

\subsection{The final induction and conclusion}

Fix $k,X$ and the parameters in \eqref{eq:precision} before applying
induction. Use the stage predicate \eqref{eq:stage-table}, with the stage
index truncated at $J$. The base and the implication to the next stage have
just been proved. Schema~\eqref{eq:qfind} therefore proves the stage-$J$
table, whose envelope is $a_J=1/q$.
Only after this finite induction do we generalize $X$. For the final target
integer $n$, take $X=n$; all subarguments used in its proof are below $n+1$.
Section~\ref{sec:pi} then yields \eqref{eq:rational-half}. Its margin over
the claimed arithmetic bound is at least $1/(2k)$, much larger than the
finite numerical error in \eqref{eq:budget-bound}. After clearing denominators,
the resulting statement is exactly \eqref{eq:arith-pnt}.

We have consequently described a derivation in the stipulated object theory,
not merely a proof in a stronger theory that the modulus is total.
There is no induction on
$\forall x\ge P_j\,\bigl(|R(x)|\le a_jx\bigr)$; that unbounded
formula was used only as readable notation before its finite-table
implementation was supplied. No assertion about arbitrary functions or
arbitrary reals enters the object-language statement.

\section{A precise route to a Lean formalization}
\label{sec:lean-plan}

\subsection{The proof-theoretic endpoint}

There are two different formalization goals.  The first is to prove the
numerical estimates in Lean's ambient mathematics.  The second, required
here, is to exhibit a derivation using only the axioms and inference rules
of intuitionistic primitive recursive arithmetic.  The first goal does not
imply the second: even a Lean proof containing no use of excluded middle
may use induction considerably stronger than quantifier-free induction.

The supplied ProveIt repository already provides a suitable deep embedding.
The following statements refer to commit
\texttt{b03925cccc76d55b9dc2a00864a5e0544b64aa7f}, inspected through the
repository tools.  The source directory is
\begin{quote}\small\ttfamily
Logic/PrimitiveRecursiveArithmetic/ConciseIPRA/\\
Lean/ConciseIPRA/
\end{quote}
Within it, \texttt{PR.lean} defines primitive-recursive function codes
\texttt{PR k}, scoped arithmetic terms \texttt{PTm n}, and the axiom family
\texttt{PRAAx}.  The function codes are generated by zero, successor,
projections, composition, and primitive recursion.  The induction axiom
requires a proof that its matrix is quantifier-free; the useful derived
rule is \texttt{PRA.qfInd}.  The module \texttt{Deduction.lean} defines
\texttt{Prf Ax Gamma phi}, intuitionistic first-order natural deduction
with equality.  Its constructors provide no excluded-middle rule.

The final theorem should therefore have the form
\begin{quote}\ttfamily
theorem pnt\_iPRA : Prf PRAAx [] pntSentence := ...
\end{quote}
Here \texttt{pntSentence} must be a \emph{defined closed formula}, with
defined primitive-recursive codes for its numerical functions.  It is not
enough to leave those codes as parameters, to prove merely that suitable
functions are primitive recursive, or to establish an asymptotic statement
about \texttt{Real.log} in Mathlib.

The existing module \texttt{NativeBiInterpretation.lean} supplies
\texttt{pra\_iff\_native}:
\[
 (\Gamma\vdash_{\mathrm{iPRA}}\varphi)
 \quad\Longleftrightarrow\quad
 \operatorname{NPrf}
   (\Gamma.\operatorname{map}\operatorname{NFm.ofPRA})
   (\operatorname{NFm.ofPRA}(\varphi)).
\]
Consequently, a completed arithmetic derivation transfers directly to the
repository's concise native calculus.  This reuses the foundational
equivalence without treating it as a source of new number-theoretic
estimates. The inspected Git blobs are recorded explicitly for reproduction:
\begin{quote}\small
\texttt{PR.lean}:\\
\texttt{3251631592070d9ed3687cc0bf7dbd5eb452e46f}\\[3pt]
\texttt{NativeBiInterpretation.lean}:\\
\texttt{ccfe499b9687ea0e176f832686d8e492fe15544b}.
\end{quote}

\subsection{An exact arithmetic statement}

For an integer $n\geq 1$, put
\[
 F(n)=(n^2)!,\qquad
 A(n)=\sum_{j=n+1}^{n^2}\frac{F(n)}{j},\qquad
 \ell(n)=\frac{A(n)}{F(n)}.
\]
Every quotient in $A(n)$ is exact.  Thus both $A$ and $F$ are functions on
natural numbers, and $\ell(n)=H_{n^2}-H_n$.  The main decidable matrix is
\[
 k\,\operatorname{dist}\bigl(\pi(n)A(n),nF(n)\bigr)\leq nF(n),
 \qquad
 \operatorname{dist}(a,b)=(a\mathbin{\dotminus}b)+(b\mathbin{\dotminus}a),
\]
where $\dotminus$ is truncated subtraction.  For $k>0$ and $n\geq 2$ this is
exactly
\[
 \left|\frac{\pi(n)\ell(n)}n-1\right|\leq\frac1k.
\]
Use the explicit modulus
\[
 N(0)=2,\qquad
 q=8k,\quad d=2^{20}q^3,\quad
 N(k)=2^{\,2^{d^2+2}}\quad(k>0).
\]
The Boolean checker returns true if $k=0$, if $n<N(k)$, or if $n<2$;
otherwise it checks the displayed integer inequality.  The closed target
is simply the universal closure of ``this checker returns true.''  It
contains no real-number objects or quantification over functions.

All its numerical constructions admit direct primitive-recursive codes.
Primality is a bounded test over $d\leq n$; $\pi$ is a running sum of its
characteristic function; factorial, products, and sums are primitive
recursions.  Division and truncated subtraction are bounded-search
constructions.  Exponentiation is primitive recursive, so the displayed
fixed nesting of exponentials causes no foundational difficulty.  Large
values affect computational practicality, not the permitted recursion
scheme.  In particular, the formalization should never attempt to expand
the closed numeral $N(1)$ during evaluation.

\subsection{Finite ceilings and arithmetic encodings}

The quantitative argument iterates an error bound.  Its implementation
must not apply induction directly to
\[
 \forall n\geq P_j\quad |\psi(n)-n|\leq a_j n,
\]
because this is not a quantifier-free induction matrix.  Instead fix an
arbitrary finite ceiling $X$ before starting the iteration.  Encode the
assertions for every relevant $n\leq X$ by a single Boolean conjunction.
For any fixed primitive-recursive Boolean predicate $b$, this conjunction
is generated by
\[
 B(0)=b(0),\qquad B(t+1)=B(t)\wedge b(t+1).
\]
Its defining equations and correctness are established by quantifier-free
induction.  The point is that $b$ is compiled as a \emph{fixed arithmetic
code}; the object language does not quantify over an arbitrary function
parameter.

For the piecewise-constant Chebyshev function, a convenient finite
invariant records both endpoint estimates.  If $U_\psi(n)$ and
$L_\psi(n)$ are rational upper and lower enclosures for $\psi(n)$, require
\[
 U_\psi(n)\leq(1+a_j)n,
 \qquad
 L_\psi(n)\geq(1-a_j)(n+1)
\]
for the relevant integers $n\leq X$.  These inequalities control the
normalized error throughout each interval $[n,n+1)$, including its left
limit at the right endpoint.  All denominators are cleared before the
Boolean test is formed.  The contraction step uses only arguments below
its current finite ceiling; its rational mesh, bounded selection of
indices, and numerical precision are also given primitive-recursive
bounds.  One then applies \texttt{PRA.qfInd} to this Boolean invariant,
and finally substitutes $X=n$.  This substitution supplies the desired
unbounded universal conclusion without using a stronger induction schema.

Finite rational enclosures are part of the proof, rather than an appeal
to numerical experiments.  Each analytic abbreviation must come with a
proved finite error estimate.  Rational sign decisions use separated
thresholds, so equality of an unknown real number to a threshold is never
tested.  The surplus in the contraction estimate pays for the enclosure
errors and the passage to endpoint bounds.  A useful implementation order
is to prove these interval contracts first, then to reuse them in every
finite summation and mesh argument.

\subsection{Module milestones and verification}

The following modules are prospective divisions of the remaining work;
their names do not claim that such completed modules already exist.
\begin{enumerate}
\item \textbf{Arithmetic codes and their defining equations.}
Compile the exact checker, including $N,A,F,\pi$, into a concrete term of
\texttt{PR 2}.  Prove that its ordinary interpretation agrees with the
executable checker.  Build reusable object-level proofs for bounded sums,
products, searches, and Boolean conjunctions.

\item \textbf{Rational intervals.}
Develop finite logarithm enclosures, exact rational arithmetic, and the
harmonic estimate for $\ell(n)$.  Every precision parameter has a specified
primitive-recursive truncation bound.  Separate these verified interval
contracts from the later conventional interpretation in the real numbers.

\item \textbf{Finite number-theoretic identities.}
Formalize factorization, the identities for the von Mangoldt and M\"obius
functions, the elementary Chebyshev bounds, and the finite summation
identities used in Selberg's formula.  Prove their rational error versions
with the constants used in the human-readable argument.

\item \textbf{The quantitative contraction.}
Formalize the bounded geometric mesh, the separated sign alternatives,
the deficit estimate, and the threshold recurrence.  Carry a finite
ceiling $X$ throughout.  The crucial outputs are explicit transformations
of the endpoint Boolean invariants, not an invocation of an unbounded
limsup argument.

\item \textbf{The modulus and the counting function.}
Verify the finite number of contraction steps and the majorization by
$N(k)$.  Complete the transfers from $\psi$ to $\vartheta$ and to $\pi$,
then replace the logarithm by $A(n)/F(n)$ with the stated error budget.

\item \textbf{The object-level certificate.}
Assemble the previous derivations into
\texttt{Prf PRAAx [] pntSentence}.  Transfer this certificate through
\texttt{pra\_iff\_native}.  Reject any remaining proof hole, additional
arithmetic axiom, or strengthened induction rule.  Record
\texttt{\#print axioms} for the exported certificate and audit the
definition of the final sentence.

\item \textbf{The semantic PNT corollary.}
Only after the arithmetic result is established, recover the conventional
statement using standard natural-number semantics and the elementary
comparison with the real logarithm.  This last theorem may be presented
using Mathlib's usual real-analysis interface; its logical status is
distinct from the object-level certificate.
\end{enumerate}

The role of the final semantic step is particularly transparent.  The
elementary inequalities for the integral of $1/t$ give
\[
 0\leq\log n-\ell(n)\leq\frac1n.
\]
Since $\pi(n)\leq n$, replacing $\ell(n)$ by $\log n$ changes the normalized
quantity by at most $1/n$.  Applying the arithmetic theorem at precision
$2k$ and taking $n\geq 2k$ yields error at most $1/k$ for
$\pi(n)\log n/n$.  Thus, for example,
\[
 \max\{N(2k),2k,2\}
\]
is a conventional modulus obtained directly from the exact arithmetic
sentence.  No hidden choice of a sufficiently large threshold occurs.

\subsection{Status of the supplied Lean artifacts}

\texttt{lean/RationalPNT.lean} contains executable definitions of the
natural-number checker and the explicit modulus, together with small
generic proofs about bounded Boolean conjunctions, prime-counting bounds,
positive factorials, and the equivalence of the finite-ceiling
specification.  It imports only \texttt{Std}.  The definition
\texttt{ArithmeticPNT} states the target; it is not a theorem proving it.

\texttt{lean/ObjectTarget.lean} supplies a precise constructor for the
closed object-language sentence and generic wrappers around the
repository's native/iPRA proof transfer.  The arithmetic code remains a
parameter in these wrappers until the compilation milestone is completed.
No PNT axiom is introduced, and no wrapper is presented as the missing
number-theoretic derivation.

\textbf{These Lean files are an outline, not a completed formalization.}
They were not compiled in the preparation environment, where no Lean,
Lake, or Elan executable was found.  The mathematical estimates and their
object-level derivations remain the formalization work described above.
The accompanying \texttt{lean/STATUS.txt} records this limitation.  A
successful future build must establish both a faithful arithmetic statement
and a derivation in the specified object theory; the absence of
\texttt{sorry} in a collection of definitions alone would not establish
either main theorem.


\section{Provenance, scope, and reproducibility}

The elementary strategy originates in Selberg's work~\cite{selberg}.
The logarithmic-square comparison with a divisor sum is discussed in
Kane's notes~\cite{kane}. Here the constants are derived explicitly,
Euler's constant is replaced by a finite harmonic parameter, and the
contraction is carried out with a rational geometric mesh and a fixed
finite arithmetic table. These details are needed for this constructive
weak-arithmetic presentation. No novelty claim about the prime number
theorem or its provability in classical elementary arithmetic is made.

The ProveIt source was inspected at commit
\begin{quote}\small\ttfamily
b03925cccc76d55b9dc2a00864a5e0544b64aa7f.
\end{quote}
The useful component for the requested Lean endpoint is the
\texttt{ConciseIPRA} project under
\texttt{Logic/PrimitiveRecursiveArithmetic/}.
The supplied Busy Beaver/ZFC and DiophantineConsistency directories were
also inspected for their scope and arithmetic infrastructure. Their
consistency theorems are not assumptions in this proof. In particular,
no consistency statement for ZFC or stronger arithmetic is used to obtain
a number-theoretic estimate. A bounded-characteristic construction found
in the DiophantineConsistency project is a useful implementation pattern,
but its classical completeness argument is not silently imported as an
intuitionistic derivation.

The archive contains the full LaTeX source, the compiled PDF, the Lean
outline, and Python programs implementing the exact arithmetic predicate
and checking finite identities and arithmetic bookkeeping. Those finite
checks do not prove the universal theorem; the mathematical proof is the
argument above. The Lean outline was not compiled because the preparation
environment had no Lean executable. A completed deeply embedded Lean
derivation remains the formalization work described in the preceding section.

\begin{thebibliography}{9}
\bibitem{selberg}
Atle Selberg, \emph{An elementary proof of the prime-number theorem},
Annals of Mathematics (2) \textbf{50} (1949), 305--313.
\url{https://www.math.lsu.edu/~mahlburg/teaching/handouts/2014-7230/Selberg-ElemPNT1949.pdf}.

\bibitem{kane}
Daniel Kane, \emph{Notes on Selberg's proof of the prime number theorem},
Winter 2025.
\url{https://cseweb.ucsd.edu/~dakane/Math104B/SelbergNotes.pdf}.

\bibitem{cd}
C. Cornaros and C. Dimitracopoulos,
\emph{The prime number theorem and fragments of PA},
Archive for Mathematical Logic \textbf{33} (1994), 265--281.
\url{https://doi.org/10.1007/BF01270626}.
The publisher's abstract was accessible; its full proof is not used here.

\bibitem{avigad}
Jeremy Avigad, \emph{Number theory and elementary arithmetic},
Philosophia Mathematica (3) \textbf{11} (2003), 257--284,
especially p.~269, footnote~14.
\url{https://www.andrew.cmu.edu/user/avigad/Papers/elementary.pdf}.

\bibitem{sudac}
Olivier Sudac, \emph{The prime number theorem is PRA-provable},
Theoretical Computer Science \textbf{257} (2001), 185--239.
\url{https://doi.org/10.1016/S0304-3975(00)00116-X}.
Cited for the logical issue explained by Avigad, not as a proof of the
intuitionistic conclusion.

\bibitem{proveit}
Vladimir Reshetnikov et al., \emph{ProveIt},
\texttt{ConciseIPRA}, inspected commit \texttt{b03925cccc76}.
\href{https://github.com/VladimirReshetnikov/ProveIt/tree/b03925cccc76d55b9dc2a00864a5e0544b64aa7f/Logic/PrimitiveRecursiveArithmetic/ConciseIPRA}{Pinned repository source}.
\end{thebibliography}
\end{document}
